\documentclass[11pt]{article}
\usepackage[margin=25mm]{geometry}
\usepackage{fontspec}
\setmainfont{Latin Modern Roman}
\usepackage{amsmath,amssymb,mathtools}
\usepackage{unicode-math}
\setmathfont{Latin Modern Math}
\usepackage{xurl,hyperref}
\hypersetup{colorlinks=true,urlcolor=blue}
\setcounter{secnumdepth}{0}
\setlength{\parindent}{0pt}
\setlength{\parskip}{5pt}
\setlength{\emergencystretch}{3em}
\newcommand{\sref}[2]{\hyperlink{source:#1:#2}{[S#2]}}
\newcommand{\cataloguescope}[1]{\paragraph{Recorded target.}#1}
\newcommand{\R}{\mathbb R}
\newcommand{\C}{\mathbb C}
\newcommand{\N}{\mathbb N}
\newcommand{\Z}{\mathbb Z}
\newcommand{\Q}{\mathbb Q}
\newcommand{\F}{\mathbb F}
\newcommand{\eps}{\varepsilon}
\begin{document}
% UnsolvedMath ID 20002522; source catalogue code AIM-PHYSICS-0030
\setcounter{section}{20002521}
\section{Why gravitational instantons arise from monopoles}\label{20002522}
{\small UnsolvedMath ID 20002522 · Mathematical Physics}
\subsection{Definitions and mathematical statement}

\paragraph{Shared notation.}$\mathbb N=\{1,2,\ldots\}$; the standard number fields have their usual meanings. All additional conventions are stated locally.


\paragraph{Definitions and notation.}A hyper-Kähler metric has three parallel complex structures satisfying the quaternionic relations. In four dimensions a gravitational instanton is a complete noncompact Ricci-flat geometry in the source’s self-dual setting. ALE and ALF describe different asymptotic geometries; ALF has an asymptotic circle direction over a three-dimensional base. A monopole moduli space parametrizes gauge-equivalence classes of solutions of a Bogomolny equation with specified charge and boundary behavior. A hyper-Kähler reduction takes a level set of three moment maps and quotients by the acting group. A triholomorphic isometry preserves all three complex structures. Wick rotation is a complex change of a time coordinate, requiring geometric conditions to yield a real Riemannian metric.

\paragraph{Catalogue wording retained for traceability.}
Retain the seven source questions: explain why gravitational instantons should arise as monopole moduli spaces; determine whether suitable eight-dimensional noncompact hyper-Kähler manifolds are moduli spaces and what happens above dimension eight; determine which instantons arise by Wick rotation beyond the Taub--NUT example; find nontrivial hyper-Kähler reductions of finite-dimensional linear spaces; decide whether every gravitational instanton has an asymptotic local triholomorphic isometry; and identify settings in which Ricci-flat Einstein solutions strictly exceed self-dual solutions. The source discusses ALF Euclidean Schwarzschild as a non-self-dual example, the ALE Ricci-flat/self-dual conjecture in dimension four, a Kähler conjecture in higher even dimensions, and a flatness argument in odd dimensions.
\paragraph{Scope and source qualification.}The workshop contains physical intuition and partial answers alongside questions. They are source context, not independent proofs checked here. Well-behaved asymptotics and the higher-dimensional replacement for self-duality require precise specification; the original informal breadth is preserved. The catalogue title/status was paired with a machine-generated research attempt; that attempt is not peer-reviewed evidence of settlement. The statement covers the original source questions rather than only the title’s subcase.
\subsection{Short English statement}
Understand moduli-space origins and constructions of gravitational instantons, their asymptotic symmetries, and distinctions among Ricci-flat and self-dual geometries.
\subsection{Research attempt}
\paragraph{Model, date and investigated scope.}
OpenAI GPT-6 Astra (configured model \texttt{gpt-6-astra}), ultra reasoning; generated 11 October 2026. Analyzed a free positive-level hyper-Kahler quotient of flat H2 and its cotangent-bundle complex geometry. The broad instanton/moduli/Wick-rotation/asymptotic/dimensional questions remain unresolved.

\paragraph{Sources and literature check.}
Read the supplied complete definition, catalogue wording, scope qualifications and source list for this canonical ID. The original-source citations are inherited as supplied references; no fresh primary-source or novelty search was performed for this attempt. No current open-status claim is inferred from this calculation.

\paragraph{A finite-dimensional hyper-Kahler quotient test.}Take \(\mathbb H^2\cong\mathbb C_z^2\oplus\mathbb C_w^2\) with its flat hyper-Kahler structure and circle action \((z,w)\mapsto(e^{i\theta}z,e^{-i\theta}w)\). In one conventional normalization the real and complex moment-map equations at positive level \(\zeta\) are
\[|z|^2-|w|^2=\zeta,\qquad z_1w_1+z_2w_2=0.
\]
The first forces \(z\ne0\), so the circle action is free on this level: fixing \(z\) forces \(e^{i\theta}=1\). The three real equations have independent differentials there. To check this, use the free orbit vector \(X\): gradients of the three moment maps are nonzero constant multiples of \(IX,JX,KX\) (depending on normalization), which are pairwise orthogonal nonzero vectors by the quaternionic relations. Thus the level has dimension five, and its circle quotient has real dimension four.
\paragraph{What geometry is visible algebraically.}The complex equation says that \(w\) annihilates the line spanned by \(z\) under the displayed bilinear pairing; quotienting by complex scaling \((z,w)\mapsto(tz,t^{-1}w)\) identifies the complex quotient with the cotangent bundle of \(\mathbb{CP}^1\). Each nonzero complex orbit has a unique positive scaling satisfying \(r^2|z|^2-r^{-2}|w|^2=\zeta\), since the left side is strictly increasing from a nonpositive limit to infinity. The remaining phase is exactly the circle quotient.
\paragraph{Gap.}This is a standard nontrivial linear-reduction model; it is not a proof that arbitrary gravitational instantons arise as monopole moduli spaces. No general completeness/asymptotic classification, Wick-rotation criterion, local triholomorphic isometry theorem, eight/higher-dimensional moduli realization or Ricci-flat versus self-dual classification is derived. Relating a specified instanton metric to a monopole moduli construction remains the first missing step in the main source question.

\paragraph{Verification and adversarial audit.}
Checked free action, moment-map rank via quaternionic gradients, dimension count, positive scaling uniqueness and scope beyond the standard quotient. This was an LLM adversarial audit, not independent human verification; no theorem prover was used. Prompts A then B applied to this ID: Checked free action, moment-map rank via quaternionic gradients, dimension count, positive scaling uniqueness and scope beyond the standard quotient. Scope remains limited: Construct and identify a monopole moduli metric for a specified instanton and verify its global geometry and asymptotics. Final pass clarified moment-map gradient normalization constants.

Coordinator checked the free positive-level circle action, moment-map rank, quotient dimensions and positive complex-orbit scaling. This additional check was model-assisted, not independent human verification.

\paragraph{Tools, process and contributions.}
Python 3 and local shell used to read supplied files and serialize this manuscript; mathematics derived symbolically in the text, with no numerical experiment or formal proof assistant. Contribution prompt A was followed by prompt B. GPT-6 Astra with ultra reasoning drafted the research and performed its stated self-audit; any additional coordinator contribution is identified in the verification and audit. A GPT-6-based Codex coordinator assembled the manuscript and checked repository integration; its exact variant and reasoning configuration were not exposed. No human mathematical contribution, independent human proof check or formal proof certification is supplied. Complete derivations and the consulted source references are included here; private reasoning transcripts are not published.

\paragraph{Final status and first remaining gap.}
\textbf{UNRESOLVED}. Construct and identify a monopole moduli metric for a specified instanton and verify its global geometry and asymptotics.

\section{Completion Estimate}
\noindent\textbf{COMPLETION: 0\%}
The calculation is an elementary benchmark or conditional reduction; it establishes no claimed new portion of the complete source target.

\subsection{Sources}
\begin{itemize}
\item[\hypertarget{source:20002522:1}{S1}] ULAM.AI and the UnsolvedMath Contributors, UnsolvedMath, record 20002522 (AIM-PHYSICS-0030). Catalogue curation under CC BY 4.0; definitions and English exposition added locally.
\url{https://huggingface.co/datasets/ulamai/UnsolvedMath}.
\item[\hypertarget{source:20002522:2}{S2}] American Institute of Mathematics, $L^2$ harmonic forms in geometry and string theory; notes by Anda Degeratu and Mark Haskins, March 2004, 17 March 2004, Self-dual Gravitational Instantons questions.
\url{https://aimath.org/WWN/l2harmonic/l2harmonic.pdf}.
\end{itemize}
\label{end:20002522}
\end{document}
